% In this file you should put the actual content of the blueprint.
% It will be used both by the web and the print version.
% It should *not* include the \begin{document}
%
% If you want to split the blueprint content into several files then
% the current file can be a simple sequence of \input. Otherwise It
% can start with a \section or \chapter for instance.

\chapter{Group transformation families}

Throughout, $G$ is a group acting on a measurable space $\Omega$ (the sample space)
and on a set $\Theta$ (the parameter space). We write $g x$ and $g \theta$ for
the actions, and $g_* \mu$ for the pushforward of a measure $\mu$ on $\Omega$
along $x \mapsto g x$.

\section{Definitions}

\begin{definition}[Group transformation family]
  \label{def:equivariant-family}
  \lean{ProbabilityTheory.IsEquivariantFamily}
  \leanok
  A family of measures $(P_\theta)_{\theta \in \Theta}$ on $\Omega$ is a
  \emph{group transformation family} for $G$ if
  \[ P_{g\theta} = g_* P_\theta \qquad \text{for all } g \in G,\ \theta \in \Theta, \]
  that is, $P_{g\theta}(A) = P_\theta(g^{-1}A)$ for every measurable $A \subseteq \Omega$.
\end{definition}

\begin{definition}[Invariant statistic]
  \label{def:invariant-statistic}
  \lean{ProbabilityTheory.IsInvariantStatistic}
  \leanok
  A statistic $T : \Omega \to \mathcal{T}$ is \emph{$G$-invariant} if
  $T(gx) = T(x)$ for all $g \in G$ and $x \in \Omega$.
\end{definition}

\begin{definition}[Maximal invariant]
  \label{def:maximal-invariant}
  \uses{def:invariant-statistic}
  An invariant statistic $T$ is a \emph{maximal invariant} if $T(x) = T(y)$ implies
  $y = gx$ for some $g \in G$; that is, $T$ separates the orbits of $G$ on $\Omega$.
\end{definition}

\begin{definition}[Equivariant estimator]
  \label{def:equivariant-estimator}
  An estimator $\delta : \Omega \to \Theta$ is \emph{equivariant} if
  $\delta(gx) = g\,\delta(x)$ for all $g \in G$ and $x \in \Omega$.
\end{definition}

\begin{definition}[Invariant loss]
  \label{def:invariant-loss}
  A loss function $L : \Theta \times \Theta \to [0, \infty]$ is \emph{invariant} if
  $L(g\theta, gd) = L(\theta, d)$ for all $g \in G$ and $\theta, d \in \Theta$.
\end{definition}

\section{Basic results}

\begin{lemma}[Invariant statistics have orbit-constant laws]
  \label{lem:invariant-law}
  \lean{ProbabilityTheory.IsEquivariantFamily.map_invariant}
  \leanok
  \uses{def:equivariant-family, def:invariant-statistic}
  Let $(P_\theta)$ be a group transformation family and $T$ a measurable invariant
  statistic, and suppose the action of $G$ on $\Omega$ is measurable. Then for all
  $g \in G$ and $\theta \in \Theta$,
  \[ T_* P_{g\theta} = T_* P_\theta. \]
\end{lemma}

\begin{proof}
  \leanok
  We have $T_* P_{g\theta} = T_* (g_* P_\theta) = (T \circ g)_* P_\theta = T_* P_\theta$,
  where the last step uses $T \circ g = T$.
\end{proof}

\begin{corollary}[Invariant statistics are ancillary under transitive actions]
  \label{cor:ancillary}
  \uses{lem:invariant-law}
  If, in addition, $G$ acts transitively on $\Theta$, then the distribution of every
  measurable invariant statistic does not depend on $\theta$; that is, $T$ is ancillary.
\end{corollary}

\begin{proof}
  \uses{lem:invariant-law}
  Given $\theta, \theta' \in \Theta$, choose $g$ with $\theta' = g\theta$ and apply
  Lemma~\ref{lem:invariant-law}.
\end{proof}

\begin{theorem}[Risk of equivariant estimators]
  \label{thm:constant-risk}
  \uses{def:equivariant-family, def:equivariant-estimator, def:invariant-loss}
  Let $(P_\theta)$ be a group transformation family, $L$ an invariant loss and $\delta$
  a measurable equivariant estimator. Then the risk
  $R(\theta, \delta) = \int L(\theta, \delta(x))\, dP_\theta(x)$ satisfies
  $R(g\theta, \delta) = R(\theta, \delta)$ for all $g$ and $\theta$. In particular, if
  $G$ acts transitively on $\Theta$, the risk of $\delta$ is constant.
\end{theorem}

\begin{proof}
  By the change of variables formula for pushforward measures,
  \begin{align*}
    R(g\theta, \delta)
      &= \int L(g\theta, \delta(x))\, d(g_* P_\theta)(x)
       = \int L(g\theta, \delta(gx))\, dP_\theta(x) \\
      &= \int L(g\theta, g\,\delta(x))\, dP_\theta(x)
       = \int L(\theta, \delta(x))\, dP_\theta(x). \qedhere
  \end{align*}
\end{proof}
